% Drafted with an AI assistant on 12 Nov 2026 from the registered design (preregistration.md and
% config/experiment.toml). Revise it in your own words. Numbers are \val{} macros.
\section{Problem Statement}
\label{sec:problem}

\textbf{Problem.} A quantum classifier is trained in a noiseless simulator and then frozen. How
much do its accuracy and the reliability of its predictions change when the same circuits run
on real IBM Heron processors on different days and devices? How much of the change is genuine
drift rather than shot noise? And which low-cost corrections recover the loss for the least QPU
time? The gap in the literature that motivates these questions is summarised in
Section~\ref{sec:gap}.

\subsection{Setting}
Let $x\in[0,\pi]^{n}$ be an input with $n=\val{n-qubits}$ features and $\rho_\theta(x)$ the state
that a fixed circuit prepares on $n$ qubits. A model reads out Pauli expectation values and
passes them to a logistic head,
\begin{align}
\hat p(y=1\mid x)&=\sigma\!\Big(w^\top\frac{\phi(x)-\mu}{s}+b\Big), \label{eq:model}\\
\phi_j(x)&=\operatorname{tr}\!\big[P_j\,\rho_\theta(x)\big], \nonumber
\end{align}
where $P_j$ is a Pauli string and $\mu$, $s$ are the mean and standard deviation of the training
features. The circuit parameters $\theta$ and the head $(w,b,\mu,s)$ are fixed after training in
simulation. On device $d$ in session $t$, a job runs the circuits with $N=\val{shots}$ shots and
returns estimates $\hat\phi_{d,t}(x)$. Their expectation $\bar\phi_{d,t}(x)$ differs from
$\phi(x)$ because of gate, readout and decoherence errors, and it changes with $t$ as the
processor drifts and is recalibrated. Shot noise adds a variance of
$(1-\bar\phi_{d,t,j}^{2})/N$ to feature $j$. Accuracy is measured on a fixed test set of
\val{n-test} inputs, giving $A_{\mathrm{sim}}$ in the simulator and $A_{d,t}$ on hardware.

\subsection{Quantities to determine}
\begin{itemize}
\item \emph{Deployment gap and emulator error.} $G_{d,t}=A_{\mathrm{sim}}-A_{d,t}$ and
$E_{d,t}=A^{\mathrm{emu}}_{d,t}-A_{d,t}$, where $A^{\mathrm{emu}}_{d,t}$ is the accuracy
predicted by a noise model built from the calibration reported at submission.
\item \emph{Drift.} For each input and feature, the ratio of the variance of the estimate across
sessions to its expected shot variance,
\begin{equation}
r_d(x,j)=\frac{\operatorname{Var}_t\big[\hat\phi_{d,t,j}(x)\big]}
{\operatorname{E}_t\big[(1-\hat\phi_{d,t,j}(x)^{2})/N\big]}.
\label{eq:drift}
\end{equation}
A ratio of one means that shot noise alone explains the variation between sessions.
\item \emph{Noise map.} The scale $\gamma$ and shift $\beta$ in
$\bar\phi_{d,t,j}(x)\approx\gamma\,\phi_j(x)+\beta$, and their stability over $t$.
\item \emph{Repairs.} A repair $\mathcal R$ maps the hardware counts, and possibly $k$ labelled
hardware inputs, to corrected predictions. Its gain is $g(\mathcal R)=A^{\mathcal R}_{d,t}-
A_{d,t}$, and its cost $c(\mathcal R)$ is the extra QPU seconds per session.
\item \emph{Uncertainty.} Prediction sets $C_\alpha(x)\subseteq\{0,1\}$ with target coverage
$1-\alpha=\val{coverage-target-frac}$, judged by the empirical coverage
$\Pr[y\in C_\alpha(x)]$ and the mean set size.
\end{itemize}

\subsection{Research questions and hypotheses}
Table~\ref{tab:rq} links each research question to the quantity above and to the hypothesis that
was registered after the pilot runs and before the first campaign session (git tag
\texttt{prereg-v1}, \val{prereg-date}). Every hypothesis has a decision rule, fixed in advance and
based on 95\% bootstrap intervals, so that its outcome does not depend on the data.

\begin{table}[t]
\centering
\caption{Research questions, the quantities measured and the registered hypotheses.}
\label{tab:rq}
\footnotesize
\setlength{\tabcolsep}{3pt}
\begin{tabular}{@{}l>{\raggedright\arraybackslash}p{2.35cm}>{\raggedright\arraybackslash}p{4.2cm}@{}}
\toprule
 & Quantity & Registered hypothesis \\
\midrule
RQ1 & $G_{d,t}$, $E_{d,t}$ & H1: the emulator overestimates hardware accuracy ($E_{d,t}>0$) \\
RQ2 & drift ratio $r_d$; regression on chain calibration & H2: the ratio exceeds one. H4: readout error is the strongest predictor of feature error \\
RQ3 & $\gamma$, $\beta$ per session & none (descriptive) \\
RQ4 & $g(\mathcal R)$, $c(\mathcal R)$, coverage & H3: a 16-label refit beats normalisation. H5: split conformal calibrated in session 1 under-covers later sessions, while ACI holds \\
\bottomrule
\end{tabular}
\end{table}

\subsection{Constraints and scope}
The experiments use the free IBM Open Plan, which gives ten minutes of processor time per 28
days. The campaign is capped at \val{campaign-cap}~s of QPU time, jobs run when a shared queue
allows, and the models, datasets and qubit chains are fixed before the first session. We study
small classifiers on classical data (\val{n-qubits} qubits, two datasets) and do not train on
hardware. We do not claim quantum advantage: the quantum models are compared with tuned
classical baselines only to put their accuracy in context. The study is a measurement of
deployment behaviour on two processors of one family, not a general model of quantum noise.
